\section{Introduction}\label{sec:introduction}

Classical Cantor sets and their thermodynamic formalism start from a fixed rule system. An iterated function system, a graph-directed construction or a subshift of finite type is given in advance, and one then studies the sets, measures and pressure functions it generates \cite{Hutchinson1981Fractals,Moran1946AdditiveFunctions,MauldinUrbanskiGDMS,FalconerFractalGeometry,Bowen1979Quasicircles,Ruelle1978ThermodynamicFormalism,WaltersErgodic1982}. This paper concerns a system whose rules are rewritten while it runs. Its state includes a $20\times 20$ Markov kernel~$K$. At every step, six coupled mechanisms read the current kernel, choose a way of looking at it (a \emph{lens}) and a way of grouping its states (a \emph{packaging}), rewrite and gate the kernel accordingly, adapt a timescale, update a budget, and add a bounded innovation. We call these mechanisms P1--P6. They are the six primitives of the Six Birds program \cite{TsiokosFoundations}. Here they are simply an explicit, deterministic-plus-innovation update law, described in Section~\ref{sec:substrate}.

Such a system can be described at two levels. The first is a \emph{growth theory}~$\Tzero$. It weights each transition matrix entrywise by a selector observable and records how the norms of the resulting ordered matrix products grow, together with the observable control data. This is the natural input to a pressure function, in the sense of products of nonnegative matrices \cite{FengLau2002MatrixPressure,CaoFengHuang2008SubAdditive}. The second is a \emph{completion theory}~$\Tone$. It applies to the current kernel an evolve--forget--reinstate map: transport one lazy step, forget the distribution within each package, and reinstate it from fixed prototypes. It then records the unique fixed point of this map. The completion theory retains everything the growth theory records and adds the fixed point. The question this paper answers is whether the addition is genuine: \emph{is the completion object determined by the growth description, or does the completion theory strictly extend it?}

A strict extension is a precise thing to establish. It requires two states with the same base description and different extended descriptions; equivalently, the extended object map does not factor through the base object map \cite{LindMarcus1995SymbolicDynamics,TsiokosFoundations}. For such a statement to mean anything, three ingredients must be explicit: the class of states, the base description and the extended object. The paper fixes all three.

\paragraph{Main results.}
We work on an explicitly constructed set~$\Sh$ of states, the \emph{controlled shell}. It contains states of the original twenty-state system together with prescribed future innovation words. The shell has two cyclic components. In each, the future kernels are chosen from twelve Cantor-parametrized families, one for each phase of a twelve-step cycle. Within either component, the admissible parameter words range over a countable product $\Cr^{\N}$ of scaled middle-thirds Cantor sets; this is the sense of ``Cantor shell'' in the title. The shell also contains finitely many entry steps that join two special states to these components. The four main results are as follows.

\begin{enumerate}[label=\textbf{\Alph*.}]
  \item \textbf{A lawful controlled shell} (Theorem~\ref{thm:shell}). The shell~$\Sh$ is nonempty and forward invariant under the unmodified update law. All kernels on it are stochastic with entries above $1/100$. Timescale and budget stay in certified ranges, the kernel keeps changing by a definite amount once per cycle, and each of the six mechanisms changes the update somewhere on~$\Sh$.
  \item \textbf{Growth pressure} (Theorem~\ref{thm:pressure}). For each of the two original selector observables, the supremum growth pressure $P_\Sh(s)$ of the weighted history products exists for every $s\ge 0$. It is Lipschitz and strictly decreasing, satisfies $P_\Sh(0)=\log 20$ and $P_\Sh(1)\le-\log(5/4)$, and has a unique zero, which lies in $(1/3,1)$.
  \item \textbf{Strict extension} (Theorem~\ref{thm:strict-extension}). There are two states $\xp,\xm\in\Sh$ whose base descriptions coincide. They agree on configuration, timescale, budget, phase, lens, reinstatement data, and the norms of their weighted history products for every horizon, every $s\ge 0$ and both observables. Yet their completion objects differ. No function of the base description recovers the completion object on~$\Sh$.
  \item \textbf{Growth blindness} (Theorem~\ref{thm:blindness}). For the same two states under any common admissible future, the full weighted history matrices coincide from horizon two onward. Their norms coincide at every horizon. Every growth rate, and in particular every pressure computed from these observations, is therefore blind to the difference between the two completion objects.
\end{enumerate}

\begin{figure}[t]
\centering
\begin{tikzpicture}[
  font=\small,
  st/.style={draw, rounded corners=2pt, minimum width=1.5cm, minimum height=0.8cm, align=center, fill=white},
  obj/.style={draw, rounded corners=2pt, minimum height=0.6cm, align=center, fill=black!5, inner sep=3pt},
  obs/.style={draw, rounded corners=2pt, align=center, fill=black!5, inner sep=4pt},
  arr/.style={-{Latex[length=2mm]}, thick},
  objarr/.style={-{Latex[length=1.6mm]}, densely dashed},
]
  % states
  \node[st] (xp) at (0, 1.0) {$\xp$\\[-2pt]\scriptsize $\Kp$};
  \node[st] (xm) at (0,-1.0) {$\xm$\\[-2pt]\scriptsize $\Km=\Kp^{\top}$};
  \node[st] (R)  at (3.0,0)  {join\\[-2pt]\scriptsize kernel $R$};
  \node[st] (y1) at (5.4,0)  {entry\\[-2pt]\scriptsize prefix};
  \node[st] (y2) at (8.1,0)  {common parameter word\\[-2pt]\scriptsize $(\theta_t)_{t\ge1}\in\Cr^{\N}$};
  \draw[arr] (xp) -- node[above, sloped, font=\scriptsize]{$N_+$} (R);
  \draw[arr] (xm) -- node[below, sloped, font=\scriptsize]{$N_-$} (R);
  \draw[arr] (R) -- (y1);
  \draw[arr] (y1) -- (y2);
  \draw[arr] (y2) -- ++(1.5,0);
  % completion objects
  \node[obj, anchor=east] (gp) at (-1.3, 1.0) {$\Gamma_{+}=\Fix E_{\xp}$};
  \node[obj, anchor=east] (gm) at (-1.3,-1.0) {$\Gamma_{-}=\Fix E_{\xm}$};
  \draw[objarr] (xp.west) -- (gp.east);
  \draw[objarr] (xm.west) -- (gm.east);
  \node[font=\scriptsize, align=center] at ($(gp.south)!0.5!(gm.north)$) {$\Gamma_{+}\neq\Gamma_{-}$};
  \node[font=\scriptsize, anchor=south] at (gp.north) {completion theory $\Tone$};
  % growth observations
  \node[font=\scriptsize, anchor=south west] at ($(R.north)+(0.1,0.35)$) {$A_s(\xp)A_s(R)=A_s(\xm)A_s(R)$};
  \node[obs, anchor=north, font=\scriptsize] (ob) at (4.2,-1.85) {growth theory $\Tzero$:\quad
     $C_{s,n}(\xp)=C_{s,n}(\xm)$ for $n\ge 2$;\quad
     $\norm{C_{s,n}(\xp)}=\norm{C_{s,n}(\xm)}$ for all $n\ge0$, $s\ge0$};
\draw[objarr] (R.south) -- (R.south |- ob.north);
\end{tikzpicture}

\caption{The split pair behind Theorems~\ref{thm:strict-extension} and~\ref{thm:blindness}. The kernels $\Kp$ and $\Km=\Kp^{\top}$ differ only by a small directed cycle on states $0\to1\to2\to0$. Their current completion objects $\Gamma_\pm$ (left) are different probability vectors. One admissible update sends both kernels to the same joining kernel~$R$. After that they follow a common future word in the Cantor-parametrized shell. Because the first weighted matrices have equal row sums and differ only in columns $0,1,2$, on whose indices the rows of $R$ coincide, the weighted history products coincide from step two onward (bottom). Every growth observation agrees, while the completion objects differ.}
\label{fig:split-merge}
\end{figure}

The mechanism behind C and D is elementary and is shown in Figure~\ref{fig:split-merge}. The two kernels are transposes of each other, and each of their rows is a rearrangement of the corresponding row of the other. Entrywise weights therefore produce equal row sums. The one-step difference lives in three columns, and the next kernel has equal rows on those three indices. Multiplying by that kernel erases the difference. The completion map, by contrast, uses the labelled off-diagonal structure of the current kernel, which the growth description never sees.

\paragraph{Consequence for pressure disintegration.}
One might expect a strict extension to have a thermodynamic signature. Disintegrate a reference law over the completion objects, and the conditional pressures should differ, giving a positive weighted gap between the global pressure and the average of the conditional ones. Theorem~\ref{thm:blindness} shows that this is not a consequence of strictness. In Section~\ref{sec:disintegration} we prove three facts:
\begin{itemize}
  \item for any fixed finite family of positive-probability fibers whose conditional pressure limits exist, the reference pressure is the largest conditional pressure, and the gap is positive exactly when two conditional pressures differ (Proposition~\ref{prop:finite-disintegration});
  \item on the same shell, with the same potential and the same two completion objects, a common future law gives an identically zero gap (Proposition~\ref{prop:zero-gap});
  \item a correlated future law on the same shell gives a certified positive gap (Proposition~\ref{prop:positive-gap}).
\end{itemize}
The gap is therefore a property of the chosen reference law, not of the structural extension.

\paragraph{How the results are certified.}
The proofs combine three kinds of argument, and Section~\ref{sec:verification} states which result rests on which. Lean~4 with mathlib \cite{deMouraUllrich2021Lean4,mathlib2020} checks the twenty-state witness, including the exact rational kernels, completion operators, stationary vectors and history identities. It also checks the abstract matrix-pressure, root-location and finite-disintegration arguments and the induction that turns one-step invariance into an all-time statement. The all-time shell inequalities, such as selector margins, innovation bounds and budget gains, are certified over boxes covering the whole shell by exact rational interval arithmetic with outward rounding \cite{MooreKearfottCloud2009Interval}. The correspondence between the implemented update formulas and their mathematical statement, the construction of the reference laws, and an imported variational theorem for subadditive cocycles \cite{Schreiber1998Subadditive} are analytic arguments, given in the text.

\paragraph{Scope.}
The shell is a constructed class. Its future innovation words are chosen to be legal for the original update. We do not claim that independently drawn random innovations keep the system in~$\Sh$ with positive probability, nor that seeded floating-point simulations reach it. All statements concern the exact-real interpretation of the update formulas. Strictness is relative to the stated base: a base that retained the full labelled kernel would determine the completion object directly. Section~\ref{sec:discussion} collects these limits.

\paragraph{Organization.}
Section~\ref{sec:substrate} describes the substrate and its update law. Section~\ref{sec:shell} constructs the controlled shell and proves Theorem~\ref{thm:shell}. Section~\ref{sec:pressure} defines the growth pressure and proves Theorem~\ref{thm:pressure}. Section~\ref{sec:completion} defines completion objects and proves that they are well defined on the whole shell. Section~\ref{sec:extension} constructs the split pair and proves Theorems~\ref{thm:strict-extension} and~\ref{thm:blindness}. Section~\ref{sec:disintegration} treats pressure disintegration. Section~\ref{sec:verification} describes the verification. Sections~\ref{sec:related-work} and~\ref{sec:discussion} discuss related work, limitations and open problems. Appendix~\ref{app:witness-data} lists the explicit witness data, and Appendix~\ref{app:version-history} records the version history.
